\section{Discussion and conclusion}
\label{sec:discussion}

The atomic source condition and directional smoothness yield an \(O(m^{-1})\) energy bound
for the ideal continuous-dictionary CGA. For finite-pool CGA, continuous-dictionary coverage
enters multiplicatively, whereas score evaluation, radial stationarity, and coefficient
optimization introduce distinct additive errors. The one-step inequality specifies the relative
accuracy conditions under which the same energy order is retained.

For standard pure, regularized, and reaction \(p\)-energies with \(p\geq2\), the energy gap
is equivalent to the squared natural distance and controls the \(p\)-th power of the Sobolev
error. An energy exponent \(\alpha\) therefore gives natural-distance and Sobolev exponents
\(\alpha/2\) and \(\alpha/p\). This conclusion uses the pointwise structure of the energy
density; a general convex density with \(p\)-growth need not satisfy the same equivalence.

The sharper estimates in \Cref{sec:c5} require additional assumptions. Residual
comparison relates the dictionary seminorm to the full dual norm. Local Hessian transfer
requires coercivity, a dictionary norm bridge, relative score resolution, and source and entropy
bounds after restart. The finite-trajectory estimates use the projections onto the computed
spaces, together with either an approximate source representation or a
fractional lower bound on the normalized innovations. The quartic energy decomposition relates
both projection estimates to nonlinear energy errors.

The concrete semilinear family in \Cref{sec:c5-hessian} makes the nonlinear interface explicit.  Its
$C_{\rm emb}$, $L_\lambda$, $C_\lambda$, and $B_\lambda$ constants give a computable candidate
margin $\Delta_\lambda$ with the required $C_\lambda^2$ factor.  The Route-A closure certificate
then requires the same projected, unrenormalized restarted dictionary to satisfy a quantitative entropy
bound and requires every declared state to remain in the Hessian ball.  These checks are now isolated
as falsifiable obligations; the manuscript does not call them an unconditional theorem before they
are verified.

The one-dimensional pure and regularized quartic calculations evaluate these quantities at
every accepted state on \(N=8\)--32 and continue the same prefixes through \(N=64\) with the
constants held fixed. Active-space Hessian spectra are positive, but the
sufficient selection margins are negative throughout this window. The continuous
Hessian-transfer hypotheses are therefore unverified. The direct projection and energy
identities agree to floating-point accuracy under two quadrature rules. Both the source and
fractional-innovation bounds cover the computed errors, although the terminal bounds remain
several thousand times larger. With the prescribed exponent \(\alpha=6\), the latter yield
finite-window comparison powers \(6\), \(3\), and \(3/2\) for energy, natural distance, and
\(W^{1,4}\) error. The fitted innovation powers are descriptive: positivity on a finite
trajectory neither determines a unique decay exponent nor gives constants uniform over
increasing widths. In the continuation, the normalized-innovation sufficient condition holds
on most transitions: 27 of 32 for the pure model and 28 of 32 for the regularized model.
The few exceptions preclude a uniform application of the frozen-constant bound across the
extended window, so the continuation remains a stepwise diagnostic.

At the largest common coefficient count, CGA has the smallest relative Sobolev error in the
four comparisons after filtering to the reported converged RFM realizations. The same ordering holds in the
two-dimensional pure \(p=4\) example, where only four of ten prescribed RFM realizations reach the terminal
width. These comparisons concern the specified implementations and coefficient counts;
they establish neither an ordering at every width nor a runtime advantage.

The remaining work is therefore specific: close the three Route-A gates for a declared nonzero $\lambda$ range,
or switch to the error-decomposition controller described in the revision plan.  More examples alone
would not close these gates.  Reliable numerical enclosures and comparisons with a common
computational-work measure would provide additional evidence beyond the finite-window comparisons
reported here.
